\begin{abstract}
We study a single frozen question: what are the formal limits of intervention
in enterprises modeled as continua under the Ontology of Continua (OC Core)
constraints~\cite{oc_core} and an executable-closed enterprise specification
(ETS.K\_EA.v1.1)~\cite{ets_kea_v11}.
Intervention is treated as an abstract operator acting on the admissible state
space $\Omega(K_{EA})$ subject to a threshold system $\Theta$, without
optimization, governance, or prescriptive semantics.

We formalize (i) feasibility as invariance under $\Theta$, (ii) boundary action
as crossing $\partial\Omega$, and (iii) outcome classes as phase labels
$\{\OK,\INERTIA,\COLLAPSE,\STOP\}$.
The main results are ``no-free'' limits: any intervention that genuinely acts
across $\partial\Omega$ and preserves all thresholds must either degrade
structural continuity $k$ or trigger a phase transition.
We conclude with explicit falsifiability criteria: the results fail if a
constructible operator can cross $\partial\Omega$, preserve $\Theta$, maintain or
increase $k$, and avoid any phase transition.
\end{abstract}
